\documentclass[11pt,a4paper]{article}

\usepackage[utf8]{inputenc}
\usepackage[T1]{fontenc}
\usepackage[margin=0.72in]{geometry}
\usepackage{amsmath,amssymb}
\usepackage{graphicx}
\usepackage{booktabs}
\usepackage{array}
\usepackage{enumitem}
\usepackage{xcolor}
\usepackage{listings}
\usepackage{hyperref}
\usepackage[capitalise,noabbrev]{cleveref}
\usepackage{float}

\hypersetup{
  colorlinks=true,
  linkcolor=black,
  citecolor=blue!60!black,
  urlcolor=blue!60!black,
}

\definecolor{codegray}{gray}{0.95}
\definecolor{codeblue}{rgb}{0.0,0.3,0.6}
\definecolor{codegreen}{rgb}{0.0,0.5,0.2}
\definecolor{codepurple}{rgb}{0.4,0.0,0.6}

\lstdefinelanguage{OCaml5}{
  keywords={let,in,rec,fun,if,then,else,match,with,begin,end,type,mutable,true,false,and,do,done,for,while,val,open,module},
  keywordstyle=\color{codeblue}\bfseries,
  comment=[s]{(*}{*)},
  commentstyle=\color{codegreen}\itshape,
  stringstyle=\color{codepurple},
  basicstyle=\ttfamily\small,
  backgroundcolor=\color{codegray},
  frame=single,
  framesep=4pt,
  breaklines=true,
  columns=fullflexible,
  keepspaces=true,
  showstringspaces=false,
}

\lstset{
  language=OCaml5,
  basicstyle=\ttfamily\small,
  backgroundcolor=\color{codegray},
  frame=single,
  framesep=4pt,
  breaklines=true,
  columns=fullflexible,
  keepspaces=true,
  showstringspaces=false,
}

\title{%
  \textbf{CS6868 Research Mini-Project} \\[0.3em]
  \Large Bounded Lock-Free Queues: From SPSC to MPMC%
}
\author{%
  Nishan Arya \\ \texttt{da25m001@smail.iitm.ac.in}
  \and
  Chirag \\ \texttt{da25m008@smail.iitm.ac.in}
  \and
  Dikshank Sihag \\ \texttt{da25m009@smail.iitm.ac.in}
}
\date{\today}

\begin{document}
\maketitle

\begin{abstract}
Lock-free bounded queues are foundational components in high-performance concurrent systems,
offering an alternative to allocation-heavy linked-list queues by operating on pre-allocated
ring buffers. This project systematically implements and evaluates three bounded queue variants
in OCaml 5---Single-Producer/Single-Consumer (SPSC), Multi-Producer/Single-Consumer (MPSC), and
Multi-Producer/Multi-Consumer (MPMC)---alongside a partial Disruptor-pipeline port.
We additionally implement the classical Michael-Scott unbounded queue and a mutex-protected
bounded queue as baselines. All implementations are verified for linearisability using
QCheck-Lin and QCheck-STM. Benchmarks across 1--8 threads reveal that: (1)~the SPSC queue
achieves up to 8.72\,M\,ops/sec with a single producer/consumer pair, delivering
competitive throughput without any CAS; (2)~the Vyukov MPMC queue and the Disruptor MPSC
pipeline both outperform the Michael-Scott queue in their respective regimes, but each
degrades sharply under high thread counts; (3)~cache-line padding via \texttt{Atomic.make\_contended}
provides a decisive improvement in contended configurations, with a 362\% gain observed
for the Disruptor MPSC at 4\,producers. We answer the central research question: the most
dramatic complexity-performance inflection occurs at the MPSC$\to$MPMC boundary, where
the benefit of per-slot sequence numbers is eroded by symmetric contention on both producer
and consumer cursors.
\end{abstract}

%==============================================================================
\section{Goals}
\label{sec:goals}
%==============================================================================

The lock-free unbounded queues studied in Lecture~08 (the Michael-Scott queue) allocate
a fresh heap node for every enqueue operation. This generates continuous pressure on the
garbage collector and causes cache misses as nodes traverse the allocator's free lists.
Bounded queues backed by pre-allocated ring buffers are the standard remedy in
performance-sensitive applications: the allocation is done once at construction time,
and thereafter every enqueue/dequeue is a pure array index operation.

The design space of bounded ring-buffer queues spans several points of increasing complexity:
\begin{itemize}[noitemsep]
  \item \textbf{SPSC} (single-producer, single-consumer): two monotonically increasing atomic indices; no CAS ever needed.
  \item \textbf{MPSC} (multi-producer, single-consumer): producers race on the tail via CAS; consumer reads head uncontested.
  \item \textbf{MPMC} (multi-producer, multi-consumer): per-slot atomic sequence numbers; both ends use CAS or fetch-and-add (FAA).
\end{itemize}
The LMAX Disruptor~\cite{disruptor} popularised ring-buffer queues in financial trading
and introduced the FAA-based multi-producer sequencer with an available-buffer bitmap to
enable out-of-order publication. This project builds all three variants in OCaml~5,
tests them exhaustively, and benchmarks them against each other and against classical alternatives.

\paragraph{Research question.}
\emph{How much throughput does a pre-allocated ring buffer gain over the
allocation-heavy Michael-Scott queue across different producer/consumer configurations,
and where in the SPSC $\to$ MPSC $\to$ MPMC design space does the complexity-performance
trade-off shift most dramatically?}

%==============================================================================
\section{Background}
\label{sec:background}
%==============================================================================

\subsection{Ring buffers and monotonic indices}

A bounded queue of capacity $C$ (a power of two) is backed by an array \texttt{buf[0..C-1]}.
Two monotonically increasing integers \texttt{head} and \texttt{tail} act as logical indices:
the slot in use at position $p$ is \texttt{buf[p \& (C-1)]}. Letting indices grow without
wrapping avoids the classic ABA ambiguity between \emph{full} ($\texttt{tail}-\texttt{head}=C$)
and \emph{empty} ($\texttt{tail}=\texttt{head}$).

\subsection{SPSC: Lamport's queue~\cite{lamport83}}

With exactly one producer and one consumer, no CAS is ever needed. The producer is the
\emph{sole writer} of \texttt{tail} and the consumer is the sole writer of \texttt{head}.
Therefore both can use plain atomic loads and stores. An atomic release-store on
\texttt{tail} after writing the value ensures the value is visible to the consumer
before the tail increment, maintaining the necessary ordering guarantee without a full
memory fence.

\subsection{MPMC: Vyukov's per-slot sequence numbers~\cite{vyukov10}}

Each slot $i$ carries its own atomic \texttt{sequence} field, initialised to $i$.
The lifecycle of a slot at logical position $p$ is:
\begin{enumerate}[noitemsep]
  \item \textbf{Free for write}: \texttt{sequence = p}. A producer who observes this and wins a CAS on \texttt{enqueue\_pos} owns the slot.
  \item \textbf{Written}: producer sets \texttt{sequence = p+1}, signalling consumers.
  \item \textbf{Free for next cycle}: consumer sets \texttt{sequence = p + C}, signalling the next producer wave.
\end{enumerate}
This requires \emph{no global lock and no ABA protection}: the sequence number encodes both the slot's generation and its state.

\subsection{Disruptor pipeline~\cite{disruptor}}

The LMAX Disruptor separates concerns differently. A \emph{sequencer} (SPSC or multi-producer)
manages a single monotonic cursor. Producers call \texttt{next} to claim a sequence number,
write to \texttt{buf[seq \& mask]}, then call \texttt{publish}. For multiple producers,
\texttt{next} uses FAA to claim a batch atomically and then writes a per-slot
\texttt{available\_buffer} flag so consumers can detect out-of-order gaps.
Consumers query a \emph{sequence barrier} which scans the available buffer to find the
highest contiguous sequence. This design is optimised for fan-out pipelines where \emph{every}
consumer processes every event, rather than work-stealing where each event goes to exactly
one consumer.

\subsection{Michael-Scott queue~\cite{ms96}}

The classical lock-free unbounded queue uses a linked list of dynamically allocated nodes.
The enqueue operation CAS-swings the tail pointer; dequeue CAS-advances the head. The
allocation of one heap node per enqueue is its primary performance liability compared to
ring-buffer designs.

%==============================================================================
\section{Tasks Undertaken}
\label{sec:tasks}
%==============================================================================

\subsection{SPSC Queue }
\label{subsec:spsc}

\paragraph{Data structure.}
The SPSC queue is a record with five fields:

\begin{lstlisting}[caption={Core SPSC record type.}]
type 'a t = {
  capacity : int;
  mask     : int;           (* capacity - 1, for cheap modulo *)
  buf      : 'a Option.t array;
  head     : int Atomic.t;  (* consumer writes; producer reads *)
  tail     : int Atomic.t;  (* producer writes; consumer reads *)
}
\end{lstlisting}

Both \texttt{head} and \texttt{tail} are allocated with
\texttt{Atomic.make\_contended}, which instructs the OCaml runtime to place
each value on its own cache line.  Without this, the two indices share a line
and every producer write invalidates the consumer's cached line and vice versa
--- this is \emph{false sharing}, discussed further in \cref{subsec:false-sharing}.

\paragraph{try\_enqueue.}
The producer reads \texttt{tail} and \texttt{head}, checks that
$\texttt{tail} - \texttt{head} < \texttt{capacity}$ (the back-pressure condition),
writes the value into \texttt{buf[tail \& mask]}, and then performs a
\emph{release store} on \texttt{tail}. The release ordering guarantees that the
value write is globally visible before the incremented tail becomes visible
to the consumer: no explicit fence is required beyond what \texttt{Atomic.set}
provides in OCaml~5's memory model.

\paragraph{try\_dequeue.}
Symmetrically, the consumer reads \texttt{head} and \texttt{tail}, returns
\texttt{None} if $\texttt{head} \geq \texttt{tail}$, reads the value, clears
the slot for the GC, and release-stores \texttt{head}. No CAS is ever issued:
the invariant that only one domain writes each index is the key correctness
property.

\begin{lstlisting}[caption={SPSC enqueue (no CAS required).}]
let try_enqueue t v =
  let tl = Atomic.get t.tail in
  let hd = Atomic.get t.head in
  if tl - hd >= t.capacity then false
  else begin
    t.buf.(tl land t.mask) <- Some v;
    Atomic.set t.tail (tl + 1); (* release store *)
    true
  end
\end{lstlisting}

\subsection{MPMC Queue : Vyukov's Algorithm}
\label{subsec:mpmc}

\paragraph{Slot structure.}
Each element of the ring buffer is not just a value but a \emph{slot} with its
own sequence number:

\begin{lstlisting}[caption={Per-slot record for Vyukov MPMC.}]
type 'a slot = {
  sequence : int Atomic.t;
  mutable  value : 'a option;
}
\end{lstlisting}

The queue-level \texttt{enqueue\_pos} and \texttt{dequeue\_pos} are plain
\texttt{Atomic.make} values (unpadded in the production implementation);
the false-sharing investigation in \cref{subsec:false-sharing} quantifies
the cost of this choice.

\paragraph{Enqueue logic.}
A producer reads \texttt{enqueue\_pos = pos}, reads \texttt{slot.sequence = seq},
and computes \texttt{diff = seq - pos}:
\begin{itemize}[noitemsep]
  \item \texttt{diff = 0}: the slot is free for this cycle; the producer CAS-es \texttt{enqueue\_pos} from \texttt{pos} to \texttt{pos+1}. If it wins, it writes the value and sets \texttt{slot.sequence = pos+1}, signalling consumers.
  \item \texttt{diff < 0}: either the buffer is full (consumer has not yet freed this slot for the current cycle) or a consumer is mid-flight; the producer checks \texttt{dequeue\_pos} to distinguish the two cases and either spins or returns \texttt{false}.
  \item \texttt{diff > 0}: another producer advanced past us; retry.
\end{itemize}

\paragraph{Dequeue logic.}
A consumer reads \texttt{dequeue\_pos = pos}, checks \texttt{slot.sequence - (pos+1) = 0}
(the slot has been written for this cycle), then CAS-es \texttt{dequeue\_pos}. On success
it reads the value and resets \texttt{slot.sequence = pos + mask + 1}, advancing the sequence
by one full buffer cycle to free the slot for the next producer wave.

An important subtlety: if \texttt{diff < 0} and \texttt{enqueue\_pos > pos}, a producer has
claimed the slot via CAS but has not yet written the value. The consumer must spin rather than
return \texttt{None}, because the enqueue has already logically committed (the producer
claimed the slot and will return \texttt{true}). Returning \texttt{None} here would violate
linearisability. The same reasoning applies symmetrically on the producer side.

\subsection{MPSC Queue }
\label{subsec:mpsc}

The MPSC queue applies Vyukov's per-slot sequence technique but distinguishes the
consumer path: because there is exactly one consumer, \texttt{dequeue\_pos} is advanced
with a plain \texttt{Atomic.set} rather than a CAS. This eliminates one atomic
read-modify-write operation on the hot consumer path.

On the producer side, the spec requires CAS (not FAA) for tail claiming. This is deliberate:
FAA irrevocably increments the counter, meaning a producer that claims a slot must always
complete the write, even if the buffer is full. CAS can fail and retry without committing,
enabling a \emph{non-blocking} \texttt{try\_enqueue} that returns \texttt{false} cleanly
when the buffer is full rather than spinning indefinitely.

A subtle addition compared to the pure Vyukov MPMC is a ``producer in-flight'' detection
on the consumer side: if the consumer finds \texttt{diff < 0} but observes that
\texttt{enqueue\_pos > pos}, it knows a producer has claimed the slot via CAS and must spin
for the value to appear, rather than incorrectly returning \texttt{None}.

\subsection{Disruptor Pipeline}
\label{subsec:disruptor}

We implement two pipeline variants:
\begin{itemize}[noitemsep]
  \item \textbf{Int\_spsc}: wraps \texttt{disruptor\_sequencer\_spsc} + \texttt{ring\_buffer} + \texttt{sequence\_barrier}.
  \item \textbf{Int\_mpsc}: wraps \texttt{disruptor\_sequencer\_mpsc} (FAA-based cursor with available-buffer bitmap).
\end{itemize}

\paragraph{SPSC Sequencer.}
The producer maintains a \emph{cached consumer sequence} to avoid reading the contended
consumer atomic on every iteration. Only when the buffer appears full does it refresh
the cache by reading the live consumer sequence. This single optimisation eliminates a
cross-domain atomic read on the common (non-full) path.

\paragraph{MPSC Sequencer.}
The multi-producer sequencer uses \texttt{fetch\_and\_add} on the cursor to claim a
batch of slots atomically:
\begin{lstlisting}[caption={FAA-based batch claim in the MPSC Disruptor sequencer.}]
let next t n =
  let hi = Sequence.fetch_and_add t.cursor n + n in
  let lo = hi - n + 1 in
  ...
  (lo, hi)
\end{lstlisting}
After writing their slots, producers mark each one in \texttt{available\_buffer} so
that out-of-order publications can be detected by the sequence barrier. A TSAN-flagged
race on \texttt{cached\_gating\_min} (a mutable field read/written by multiple producer
domains) was fixed by changing it to an \texttt{Atomic.t}: the cached value is only a
performance hint, so any recent valid observation is safe regardless of which producer
writes it last.

\paragraph{Why no Disruptor MPMC.}
FAA is irrevocable: a consumer that increments \texttt{dequeue\_pos} must complete its
read. This makes a non-blocking \texttt{try\_dequeue} interface impossible without a
serialising mutex, which defeats the purpose of FAA-based concurrency.
Furthermore, clean shutdown in an \texttt{nP:nC} sweep is inherently difficult:
a consumer that has claimed a slot via FAA but finds no value published (because the
producer was preempted between its FAA and its \texttt{available\_buffer} write) will
spin indefinitely. We therefore test Disruptor only in its natural SPSC and MPSC
pipeline topologies.

\subsection{Supporting Modules}

\texttt{sequence.ml} provides a \texttt{Sequence.t} --- a padded atomic integer
(\texttt{Atomic.make\_contended}) with get, set, cas, and fetch\_and\_add operations.
\texttt{sequence\_barrier.ml} implements the consumer-side wait: it polls the sequencer's
cursor (and, for MPSC, scans the available buffer) to find the highest contiguous
sequence that can be safely read. \texttt{wait\_strategy.ml} provides the
\texttt{Domain.cpu\_relax} spinning strategy used throughout.

%==============================================================================
\section{Testing and Verification}
\label{sec:testing}
%==============================================================================

\subsection{Smoke Test}

A hand-written smoke test (\texttt{test\_spsc\_smoke.ml}) verifies basic
enqueue/dequeue ordering, full-buffer rejection, and empty-queue return values
for the SPSC queue under sequential execution.

\subsection{Stress Tests}

Three concurrent stress tests verify that no items are lost or duplicated:
\begin{itemize}[noitemsep]
  \item \textbf{SPSC stress}: 1\,M items sent from one producer domain to one consumer domain; consumer verifies every item is received exactly once.
  \item \textbf{MPSC stress}: 4 producers $\times$ 250\,K items; consumer collects all 1\,M items and checks the union.
  \item \textbf{MPMC stress}: 4 producers $\times$ 4 consumers $\times$ 250\,K items; a shared counter verifies total items consumed.
\end{itemize}
All three pass deterministically, confirming no data races on the value path.

\subsection{QCheck-STM and QCheck-Lin }

The most rigorous testing uses \texttt{QCheck-STM} and \texttt{QCheck-Lin} from
the \texttt{multicoretests} library. We model all three queues against the same
sequential specification:

\begin{lstlisting}[caption={Sequential model shared by all three STM/Lin test suites.}]
type state = { contents : int list; capacity : int }
type cmd   = Enqueue of int | Dequeue

let next_state cmd st = match cmd with
  | Enqueue v ->
    if List.length st.contents >= st.capacity then st
    else { st with contents = st.contents @ [v] }
  | Dequeue ->
    (match st.contents with
     | []     -> st
     | _ :: t -> { st with contents = t })
\end{lstlisting}

Because SPSC is inherently single-producer/single-consumer, direct multi-domain
testing by the STM harness would violate the queue's precondition. We wrap each
side in a per-side mutex in the test adapter, so the underlying queue code is
exercised with correct single-writer invariants while the harness can still
exercise interleaved concurrent command sequences.

Three test categories were run for each of the three queue implementations:
\begin{enumerate}[noitemsep]
  \item \textbf{STM sequential} (\texttt{agree\_test}): ensures single-threaded behaviour matches the sequential model.
  \item \textbf{STM parallel} (\texttt{agree\_test\_par}): generates random parallel command sequences and checks that some sequential interleaving of the observed results is consistent with the model.
  \item \textbf{Lin linearisability} (\texttt{lin\_test}): independently verifies linearisability using the QCheck-Lin checker.
\end{enumerate}
All 9 tests (3 queues $\times$ 3 test types) pass.

\subsection{Comparison Suite }

\texttt{test\_comparison.ml} runs the same STM/Lin battery head-to-head between the
spec-correct queues and the Disruptor pipeline variants, plus a deterministic replay
test that enqueues a fixed sequence and checks dequeue order. 11 tests pass across
SPSC and MPSC design points.

%==============================================================================
\section{Evaluation}
\label{sec:evaluation}
%==============================================================================

\subsection{Experimental Setup}

All benchmarks were run on a multi-core Linux machine with OCaml~5. Each measurement
runs for 2 seconds with a fixed ring buffer of 65\,536 slots. Throughput is reported
in millions of operations per second (M\,ops/sec); higher is better. No warm-up
period is subtracted: the 2-second window is long enough to dominate start-up noise.

\subsection{Thread Sweep }
\label{subsec:thread-sweep}

\Cref{tab:1p1c,tab:npnc,tab:np1c,tab:1pnc} report throughput across four
producer/consumer configurations.

\paragraph{1P:1C baseline (\cref{tab:1p1c}).}
With a single producer and consumer, all implementations are on an even footing:
there is no contention on either index. The MPSC queue achieves the highest
throughput (24.62\,M\,ops/sec), which is counterintuitive at first glance. The
reason is that the MPSC queue uses per-slot sequence numbers (cheaper than
Option-boxing) and benefits from the fact that its dequeue path needs no CAS. The
SPSC queue (6.83\,M) is competitive, the Disruptor SPSC (8.56\,M) benefits from its
cached-consumer optimisation, and the Vyukov MPMC (10.03\,M) beats the Michael-Scott
queue (6.72\,M) despite its extra sequence-number reads. The mutex-protected bounded
queue (4.66\,M) is the slowest, as expected.

\begin{table}[H]
\centering
\caption{1P:1C throughput (M ops/sec). Single producer, single consumer.}
\label{tab:1p1c}
\begin{tabular}{@{}lc@{}}
\toprule
Implementation & M ops/sec \\
\midrule
MPSC queue (Vyukov)         & 24.62 \\
Vyukov MPMC                 & 10.03 \\
Disruptor SPSC              & 8.56  \\
SPSC queue                  & 6.83  \\
Michael-Scott               & 6.72  \\
Disruptor MPSC              & 6.30  \\
Bounded (mutex)             & 4.66  \\
\bottomrule
\end{tabular}
\end{table}

\paragraph{nP:1C (\cref{tab:np1c}).}
This is the MPSC regime. The Disruptor MPSC dramatically outperforms all others
at high thread counts, reaching 12.59\,M\,ops/sec at 6 producers. This is the
key advantage of FAA-based claiming: each producer's claim is unconditional, so
there is no CAS retry storm. However, the MPSC queue collapses from 25.83\,M
at 1 producer to 2.00\,M at 5+ producers, suffering from CAS contention. The
Vyukov MPMC and Michael-Scott queues show similar degradation curves.

\begin{table}[H]
\centering
\caption{nP:1C throughput (M ops/sec). n producers, 1 consumer.}
\label{tab:np1c}
\begin{tabular}{@{}lccccccccc@{}}
\toprule
Impl $\backslash$ n & 1 & 2 & 3 & 4 & 5 & 6 & 7 & 8 \\
\midrule
MPSC queue      & 25.83 & 5.23 & 3.25 & 2.38 & 2.00 & 2.00 & 1.91 & 2.00 \\
Disruptor MPSC  & 7.10  & 7.40 & 9.69 &10.21 &10.10 &12.59 &12.42 &12.54 \\
Vyukov MPMC     & 7.83  & 3.57 & 2.53 & 2.35 & 2.16 & 1.87 & 1.76 & 1.22 \\
Michael-Scott   & 5.96  & 3.76 & 2.98 & 2.47 & 2.15 & 2.04 & 1.79 & 1.81 \\
Bounded(mutex)  & 3.88  & 3.32 & 2.23 & 2.04 & 1.70 & 1.66 & 1.41 & 1.33 \\
\bottomrule
\end{tabular}
\end{table}

\paragraph{1P:nC and nP:nC (\cref{tab:1pnc,tab:npnc}).}
In these configurations only Vyukov MPMC, Michael-Scott, and Bounded(mutex) apply;
SPSC and MPSC variants are excluded as they do not support multiple consumers. All
three degrade roughly proportionally with thread count: adding a second consumer
roughly halves throughput, reflecting the cost of CAS contention on the consumer
side. Bounded(mutex) is surprisingly competitive in the 1P:nC case: the single
consumer mutex serialises dequeues, but the fixed cost per operation is low. The
Vyukov MPMC consistently outperforms Michael-Scott by roughly 1.2$\times$--1.5$\times$
across these configurations.

\begin{table}[H]
\centering
\caption{1P:nC throughput (M ops/sec). 1 producer, n consumers.}
\label{tab:1pnc}
\begin{tabular}{@{}lcccccccc@{}}
\toprule
Impl $\backslash$ n & 1 & 2 & 3 & 4 & 5 & 6 & 7 & 8 \\
\midrule
Vyukov MPMC    & 5.76 & 3.48 & 2.92 & 2.61 & 2.26 & 1.54 & 1.26 & 1.19 \\
Michael-Scott  & 6.08 & 3.02 & 2.24 & 1.81 & 1.42 & 1.22 & 1.08 & 1.02 \\
Bounded(mutex) & 4.21 & 3.30 & 3.02 & 2.35 & 2.07 & 1.89 & 1.94 & 1.63 \\
\bottomrule
\end{tabular}
\end{table}

\begin{table}[H]
\centering
\caption{nP:nC throughput (M ops/sec). n producers, n consumers.}
\label{tab:npnc}
\begin{tabular}{@{}lcccccccc@{}}
\toprule
Impl $\backslash$ n & 1 & 2 & 3 & 4 & 5 & 6 & 7 & 8 \\
\midrule
Vyukov MPMC    & 6.92 & 2.97 & 1.83 & 1.67 & 1.84 & 1.46 & 1.28 & 1.21 \\
Michael-Scott  & 5.74 & 2.74 & 1.95 & 1.62 & 1.43 & 1.29 & 1.22 & 1.14 \\
Bounded(mutex) & 3.87 & 3.67 & 2.28 & 2.19 & 2.06 & 1.77 & 1.69 & 1.58 \\
\bottomrule
\end{tabular}
\end{table}

\subsection{Batch-Size Sweep }
\label{subsec:batch}

\Cref{tab:batch} shows the effect of batch size $k$ on throughput for the 1P:1C configuration.

\begin{table}[H]
\centering
\caption{Batch-size sweep — 1P:1C (M ops/sec). Higher $k$ = more items per loop turn.}
\label{tab:batch}
\begin{tabular}{@{}lccccc@{}}
\toprule
Impl $\backslash$ batch $k$ & 1 & 4 & 16 & 64 & 256 \\
\midrule
Disruptor SPSC  & 6.09 & \textbf{32.04} & 7.30 & 7.31 & 7.71 \\
Vyukov MPMC     & 6.62 & 8.17  & 7.61 & 8.53 & \textbf{20.71} \\
Disruptor MPSC  & 6.79 & 6.82  & 6.37 & 6.55 & 6.44 \\
Michael-Scott   & 5.08 & 2.74  & 4.29 & 5.43 & 6.34 \\
Bounded(mutex)  & 4.00 & 3.62  & 3.95 & 3.71 & 3.69 \\
\bottomrule
\end{tabular}
\end{table}

The Disruptor SPSC shows a dramatic 5$\times$ boost at $k=4$ (32.04\,M\,ops/sec)
because the consumer drains all available events per barrier pass, amortising the
barrier-scan cost across multiple items. The effect diminishes at $k=16$ and beyond
as the per-item cost re-dominates. The Vyukov MPMC shows a late spike at $k=256$
(20.71\,M), consistent with its CAS-retry cost being amortised over a large batch.
The Disruptor MPSC is relatively insensitive to batch size, suggesting its bottleneck
is the FAA latency rather than per-item overhead.

\subsection{False Sharing }
\label{subsec:false-sharing}

\Cref{tab:false-sharing} quantifies the throughput impact of cache-line padding
on atomic indices. Padding is applied via \texttt{Atomic.make\_contended}; the
unpadded baseline uses \texttt{Atomic.make}.

\begin{table}[H]
\centering
\caption{False-sharing investigation. Delta = (padded $-$ unpadded) / unpadded $\times$ 100\%.}
\label{tab:false-sharing}
\begin{tabular}{@{}llccc@{}}
\toprule
Implementation & Config & Padded (M/s) & Unpadded (M/s) & Delta \\
\midrule
Disruptor SPSC & 1P:1C  & 7.20 & 7.36 & $-2.2\%$   \\
Disruptor MPSC & 1P:1C  & 5.56 & 4.97 & $+11.9\%$  \\
Disruptor MPSC & 4P:1C  & 8.37 & 1.81 & $+362.7\%$ \\
Vyukov MPMC    & 1P:1C  & 5.80 & 7.54 & $-23.1\%$  \\
Vyukov MPMC    & 2P:1C  & 3.83 & 2.73 & $+40.1\%$  \\
Vyukov MPMC    & 4P:4C  & 2.24 & 1.40 & $+59.9\%$  \\
\bottomrule
\end{tabular}
\end{table}

The results reveal two distinct patterns:

\paragraph{Padding helps under contention.} The 362.7\% gain for Disruptor MPSC at 4P:1C
is the headline result. With 4 producer domains all writing to an \texttt{available\_buffer}
indexed near the cursor, every cursor write invalidates the cache line shared by
other producers, causing a cache-coherence storm. Padding the cursor sequence to its own
line eliminates this entirely.

\paragraph{Padding can hurt in uncontended cases.} Vyukov MPMC at 1P:1C shows a 23.1\%
\emph{reduction} with padding. When only one producer and one consumer access each index,
padding spreads them across more cache lines, reducing spatial locality and increasing
the number of TLB entries consumed. The breakeven point is reached at 2P:1C (+40.1\%
improvement with padding), where inter-domain false sharing begins to dominate.

\paragraph{Implications for production code.} The SPSC queue uses \texttt{Atomic.make\_contended}
for both \texttt{head} and \texttt{tail} unconditionally, which is correct: in any realistic
deployment the producer and consumer run on different cores, and the $-2.2\%$ overhead
at 1P:1C is negligible compared to the protection afforded. The Vyukov MPMC production
code uses \texttt{Atomic.make} (unpadded) for its queue-level indices, which is
suboptimal at higher thread counts; \texttt{make\_contended} should be used for any
deployment with 2+ producers or 2+ consumers.

%==============================================================================
\section{Reflection on the Use of LLMs}
\label{sec:llm}
%==============================================================================

\paragraph{Tools and models used.}
We used Claude (Sonnet 4.6) for code generation assistance, debugging, and report drafting.
GitHub Copilot was used inline within the editor for boilerplate completion. ChatGPT (GPT-4o)
was consulted for explaining certain aspects of the OCaml 5 memory model and for
cross-checking our understanding of the Vyukov queue algorithm.

\paragraph{What worked well.}
Claude was highly effective at generating the QCheck-STM adapter boilerplate, which is
structurally repetitive across the three queue variants. Given a description of the
sequential model and the queue interface, it produced a correct \texttt{module SpscAdapter}
on the first attempt, including the subtle choice to wrap each side in a per-side mutex
for the test harness without modifying the underlying queue. It also correctly explained
the slot-lifecycle invariants in Vyukov's algorithm when we described the \texttt{diff}
cases in informal terms.

\paragraph{What did not work.}
The LLM initially generated an incorrect \texttt{available\_buffer} indexing formula
for the MPSC Disruptor sequencer (using \texttt{pos mod size} instead of
\texttt{(pos lsr index\_shift) land mask}), which produced a subtle bug where
out-of-order detection failed for large sequence numbers. TSAN did not catch this
because it is a logical error, not a data race; we caught it by running the STM
parallel tests which found a failing linearisability trace.

\paragraph{What was surprising.}
Both positively and negatively: the LLM was better than expected at explaining
\emph{why} the sequence-barrier spin is necessary (the ``producer in-flight'' case),
but significantly worse than expected at reasoning about the precise memory ordering
required. When asked ``does this need a fence?'' it consistently answered with
hedges like ``it depends on the architecture'' rather than reasoning from OCaml~5's
defined memory model.

\paragraph{What was difficult.}
Justifying the linearisation points for the MPSC dequeue ``producer in-flight'' spin
required genuine understanding of what linearisability means: the linearisation point
of a successful \texttt{try\_enqueue} is the CAS on \texttt{enqueue\_pos}, not the
\texttt{Atomic.set slot.sequence} that publishes the value. Once the CAS commits,
the enqueue is logically complete; any subsequent \texttt{try\_dequeue} at that slot
must eventually return the value, not \texttt{None}. The LLM could not reason about
this precisely and kept conflating ``the value is visible'' with ``the operation is
complete.'' This was the part of the project we had to reason through entirely ourselves.

\paragraph{Overall assessment.}
The LLM net-helped on scaffolding, boilerplate, and initial drafts, but was unreliable
on the subtle concurrent-correctness reasoning that constitutes the core intellectual
contribution of the project. We would use it similarly again, but with more scepticism
on anything touching memory ordering or linearisation arguments.

%==============================================================================
\section{Conclusions}
\label{sec:conclusions}
%==============================================================================

\paragraph{Answering the research question.}
A pre-allocated ring buffer gains substantial throughput over the Michael-Scott queue
in almost every configuration: 1.3$\times$--3.7$\times$ at 1P:1C depending on the
specific variant. The advantage is largest in the SPSC and MPSC regimes where per-slot
contention is absent or one-sided. The advantage narrows in the MPMC regime: at 4P:4C,
the Vyukov MPMC leads the Michael-Scott queue by only $\sim$1.03$\times$ (1.67 vs
1.62\,M\,ops/sec), because the CAS contention on both ends largely erases the
allocation savings.

\paragraph{The MPSC$\to$MPMC inflection.}
The most dramatic complexity-performance inflection occurs precisely at the MPSC$\to$MPMC
boundary. Moving from MPSC (CAS on tail only) to MPMC (CAS on both tail and head) doubles
the CAS burden under load. Meanwhile the Disruptor's FAA-based MPSC grows with more producers
rather than shrinking, reaching 12.59\,M\,ops/sec at 6 producers --- a fundamentally different
scaling curve that arises because FAA does not involve a retry loop.

\paragraph{False sharing is the dominant micro-architectural effect.}
The 362\% throughput gain from padding the Disruptor MPSC cursor is the single most
surprising result. It confirms that for shared-memory concurrent data structures on
multi-core CPUs, cache-line discipline is not an optional optimisation but a
correctness-adjacent concern: an unpadded implementation can be 4$\times$ slower
than the padded one under realistic thread counts.

\paragraph{Limitations and future work.}
Our benchmarks were conducted on a single machine; results on machines with
NUMA topology or different cache-line sizes (e.g., Apple Silicon) may differ.
The Disruptor MPMC was not implemented due to the FAA irrevocability problem;
a future implementation could use a hybrid approach (FAA with a cancellation
protocol) to recover non-blocking semantics. It would also be valuable to study
the queues under realistic workloads with non-trivial item-processing costs,
where the relative weight of queue overhead vs.\ application work changes the
optimal design point.

%==============================================================================
\section*{Contributions}
%==============================================================================

\begin{center}
\begin{tabular}{@{}lcl@{}}
\toprule
\textbf{Member} & \textbf{Contribution (\%)} & \textbf{Main responsibilities} \\
\midrule
Nishan   & 33\% & SPSC, MPSC implementations; QCheck-Lin/STM tests; stress tests \\
Chirag   & 34\% & MPMC and Disruptor pipeline; false-sharing benchmark, presentation \\
Dikshank & 33\% & Thread/batch sweep benchmarks; evaluation analysis; report \\
\midrule
\textbf{Total} & \textbf{100\%} & \\
\bottomrule
\end{tabular}
\end{center}

%==============================================================================
\bibliographystyle{plain}
\bibliography{references}
%==============================================================================

\end{document}